% ============================================================
%  Altar Valley Berms — Figure Registry Report  (Paper 1)
%  Paper 1: Berm condition and vegetation response
%  Auto-generated from figures/paper1/figure_registry_concise.txt
%  Date: 2026-02-28
% ============================================================
\documentclass[11pt, a4paper]{article}

\usepackage[utf8]{inputenc}
\usepackage[T1]{fontenc}
\usepackage{lmodern}
\usepackage[left=2.5cm, right=2.5cm, top=2.5cm, bottom=2.5cm]{geometry}
\usepackage{graphicx}
\usepackage{tabularx}
\usepackage{booktabs}
\usepackage{caption}
\usepackage{float}
\usepackage{xcolor}
\usepackage{parskip}
\usepackage{microtype}
\usepackage{hyperref}
\hypersetup{colorlinks=true, linkcolor=black, urlcolor=blue}
\usepackage{csvsimple}
\usepackage{longtable}

% ── column types ─────────────────────────────────────────────
\newcolumntype{L}{>{\raggedright\arraybackslash}X}   % wrapping left col

% ── helper: shaded row label ─────────────────────────────────
\newcommand{\rowlbl}[1]{\textbf{\small #1}}

% ── figure + registry table macro ────────────────────────────
%   \figentry{label}{file}{updated}{stats text}{interp text}
\newcommand{\figentry}[5]{%
  \subsection*{#1}
  \begin{figure}[H]
    \centering
    \includegraphics[width=\linewidth]{../figures/paper1/#2}
    \caption*{\textit{File:} \texttt{#2} \quad \textit{Updated:} #3}
  \end{figure}
  \vspace{0.5em}
  \begin{tabularx}{\linewidth}{@{} p{3.2cm} L @{}}
    \toprule
    \rowlbl{Field} & \rowlbl{Content} \\
    \midrule
    \rowlbl{Figure} & #1 \\[2pt]
    \rowlbl{File} & \texttt{\small #2} \\[2pt]
    \rowlbl{Updated} & #3 \\[6pt]
    \rowlbl{Statistical test} & #4 \\[6pt]
    \rowlbl{Interpretation} & #5 \\
    \bottomrule
  \end{tabularx}
  \vspace{3em}
}

% ============================================================
\begin{document}

\begin{center}
  {\LARGE \textbf{Geomorphic and Soil Controls on Berm Structural Integrity and Vegetation Response in the Altar Valley, Arizona}} \\[0.3em]
  {\large Figure Registry Report} \\[0.3em]
  {\small Generated 2026-02-28}
\end{center}

\vspace{1em}
\hrule
\vspace{1.5em}

This document summarises each registered figure for Paper 1: the figure image is
shown above a two-column table giving the filename, last-updated date, the
statistical test applied, and a brief interpretation of the key result.

\vspace{1em}

% ── Fig 1 ────────────────────────────────────────────────────
%% FIG_1_START
\figentry%
  {fig1}%
  {fig1_study_area.png}%
  {2026-06-06 08:06}%
  {Study area map showing 775 berms across three landform types in the Altar Valley, southern Arizona. Fan terraces: 276, Stream terraces: 140, Flood plains: 359. Panel A: location within Arizona (satellite). Panel B: berm locations overlaid on landform polygons (satellite). Berms not intersecting any mapped landform polygon assigned to Fan terraces.}%
  {Berms are distributed across the three principal landform types in the Altar Valley. Fan terraces and flood plains contain the majority of berms, with stream terraces forming a narrower transitional band.}
%% FIG_1_END

% ── Fig 2 ────────────────────────────────────────────────────
%% FIG_2_START
\figentry%
  {fig2}%
  {fig2_veg_response.pdf}%
  {2026-04-08}%
  {Vegetation response metric illustration for an example berm. Panel~A: upslope, downslope, and background area delineation. Panel~B: SAVI values in these three areas. Panel~C: monthly climatology of SAVI showing the upslope--downslope contrast used to compute the percent-difference metric $\Delta S$.}%
  {Panel A shows the spatial delineation of upslope, downslope, and background areas around an example berm. Panels B and C illustrate the SAVI time series and climatology that underpin the $\Delta S$ vegetation response metric.}
%% FIG_2_END


% ── SI: Drop-column importance ───────────────────────────────
%% SI_DROP_COLUMN_IMPORTANCE_START
\figentry%
  {SI_drop_column_importance}%
  {SI_drop_column_importance.png}%
  {2026-03-02 18:44}%
  {Drop-column (leave-one-out) importance: Random Forest (n_estimators=300, class_weight='balanced'), 5-fold stratified CV. Importance = decrease in mean CV AUC when predictor is removed from the full model. Top 5 predictors shown per outcome. Full-model CV AUC: condition = 0.73 ± 0.01, vegetation response = 0.65 ± 0.04.}%
  {Drop-column importances corroborate permutation importance rankings. Removing flow accumulation or berm length produces the largest AUC drop for condition; soil properties and landform show the largest effect for vegetation response.}
%% SI_DROP_COLUMN_IMPORTANCE_END

% ── Fig 5 ────────────────────────────────────────────────────
%% FIG_5_START
\figentry%
  {fig5}%
  {fig5_model_performance_importance.png}%
  {2026-06-06 08:07}%
  {Random Forest (n_estimators=300, class_weight='balanced'), 5-fold stratified CV. Permutation importance = mean decrease in holdout AUC from shuffling each predictor (30 repeats). Importances aggregated across one-hot columns for categorical predictors. CV AUC (condition) = 0.71 ± 0.03 (n=775). CV AUC (vegetation response) = 0.64 ± 0.05 (n=775).}%
  {Panel A: the predictor rankings differ between outcomes — flow accumulation and berm length dominate condition prediction, while hillslope gradient and soil properties are more important for vegetation response. CV AUC: condition = 0.71, vegetation response = 0.64.}
%% FIG_5_END

% ── Fig 7 ────────────────────────────────────────────────────
%% FIG_7_START
\figentry%
  {fig7}%
  {fig7_vegetation_response.png}%
  {2026-06-06 08:07}%
  {2x3 panel figure. Row 1: berm length, slope class, landform. Row 2: soil texture, B horizon presence, flow accumulation. Two-category panels (length, slope, B horizon, flow accumulation) use two-sided Fisher's exact tests. Multi-category panels (landform, texture) use pairwise two-sided Fisher's exact tests with Benjamini-Hochberg FDR correction. Flow accumulation (30 m) binarized at 2k contributing cells (two-sided Fisher p = 0.0001139).}%
  {Six predictors of vegetation response shown in a 2x3 grid. Row 1: berm length (threshold 60 m), slope class, and landform. Row 2: soil texture (n > 100 per class), B horizon presence, and flow accumulation. Vegetation response is significantly associated with slope class, landform, and soil texture. Steep-slope berms and flood-plain berms show the highest rates of vegetation response.}
%% FIG_7_END

% ── SI Fig 1 ─────────────────────────────────────────────────
%% SI_FIG_1_START
\figentry%
  {Supplementary Figure 1}%
  {SI\_fig1\_fa30\_pointplot\_by\_landform.png}%
  {2026-03-13 05:58}%
  {Point plots show mean $\pm$ 95\% CI of 30\,m flow accumulation (log scale) by landform, grouped by vegetation response (Panel A) and structural condition (Panel B).}%
  {Flood-plain berms experience the highest contributing flow accumulation regardless of outcome, while fan-terrace berms receive the least. Within each landform, effective berms tend to have similar or higher flow accumulation than ineffective berms, and degraded berms tend toward higher values than intact berms, consistent with greater hydraulic stress on degraded structures.}
%% SI_FIG_1_END

% ── SI Fig 5 ─────────────────────────────────────────────────
%% SI_FIG_5_START
\figentry%
  {Supplementary Figure 2}%
  {SI\_landform\_by\_predictors.png}%
  {2026-02-27 08:43}%
  {Two-sided binomial test ($H_0$: 50/50 split) applied within each landform
   group for each predictor category (length, slope, clay content, soil
   development); labels: *** $p < 0.001$, ** $p < 0.01$, * $p < 0.05$,
   ns~$p \geq 0.05$.}%
  {Predictors are non-randomly distributed across landforms --- high-clay soils
   concentrate on flood plains, steep slopes on fan terraces, and B-horizon
   development on fan and stream terraces --- confirming that landform is a
   meaningful stratification variable and motivating its inclusion as a covariate.}
%% SI_FIG_5_END

\newpage
\section*{Supplementary Tables}

\subsection*{Supplementary Table 1: Univariate GLM Rankings for Vegetation Response Predictors}

These tables present univariate predictor rankings using binomial generalized linear models (GLMs) to identify the most informative individual predictors for each outcome. For each predictor, we report:
\begin{itemize}
  \item \textbf{McFadden pseudo-R\textsuperscript{2}}: measures improvement in log-likelihood relative to the null model (higher values indicate better fit)
  \item \textbf{Tjur R\textsuperscript{2}}: measures separation in predicted probabilities between outcome classes
  \item \textbf{CV AUC}: 5-fold cross-validated area under the ROC curve (higher values indicate better predictive performance)
  \item \textbf{LRT p-value}: likelihood ratio test comparing the model to the null model
\end{itemize}

These univariate rankings complement the multivariate Random Forest analysis (Figure 1) by showing how each predictor performs in isolation, helping to identify the strongest individual associations before accounting for interactions and non-linear effects.

\vspace{1em}

\begin{table}[H]
\centering
\caption*{\textbf{Top 15 predictors for Vegetation Response}}
\small
\csvreader[
  tabular=lrrrr,
  table head=\toprule \textbf{Predictor} & \textbf{McFadden $R^2$} & \textbf{Tjur $R^2$} & \textbf{CV AUC} & \textbf{LRT $p$} \\ \midrule,
  late after line=\\,
  table foot=\bottomrule
]{../figures/paper1/SI_table_predictors_vegetation_response_clean.csv}
{predictor=\pred,mcfadden_r2=\mcfad,tjur_r2=\tjur,cv_auc=\cvauc,lrt_p=\lrtp}
{\pred & \mcfad & \tjur & \cvauc & \lrtp}
\end{table}

\vspace{2em}

\subsection*{Supplementary Table 2: Univariate GLM Rankings for Structural Integrity Predictors}

\begin{table}[H]
\centering
\caption*{\textbf{Top 15 predictors for Structural Integrity}}
\small
\csvreader[
  tabular=lrrrr,
  table head=\toprule \textbf{Predictor} & \textbf{McFadden $R^2$} & \textbf{Tjur $R^2$} & \textbf{CV AUC} & \textbf{LRT $p$} \\ \midrule,
  late after line=\\,
  table foot=\bottomrule
]{../figures/paper1/SI_table_predictors_structural_integrity_clean.csv}
{predictor=\pred,mcfadden_r2=\mcfad,tjur_r2=\tjur,cv_auc=\cvauc,lrt_p=\lrtp}
{\pred & \mcfad & \tjur & \cvauc & \lrtp}
\end{table}

%% SI_FIG_4_START
\figentry%
  {Supplementary Figure 3}%
  {SI\_fig3\_countplots\_by\_outcome\_and\_landform.png}%
  {2026-03-13 07:33}%
  {3$\times$3 panel of count plots. Rows: soil development (B horizon vs.\ none), berm length class (short vs.\ long), slope class (shallow vs.\ steep). Columns: vegetation response, structural condition, landform.}%
  {Fan-terrace berms are overwhelmingly found in B-horizon soils, whereas stream-terrace and flood-plain berms dominate in soils without a B horizon. Berm length and slope class are more evenly distributed across outcomes, though short berms are disproportionately intact and steep-slope berms show a higher proportion of vegetation response.}
%% SI_FIG_4_END


%% fig_rf_comparison_START
\figentry%
  {RF Full vs Binary}%
  {fig_rf_full_vs_binary_comparison.png}%
  {2026-03-08 09:54}%
  {Random Forest feature importance comparison: Full vs Binary predictor sets for Condition (Intact) and Vegetation Response (Effective) outcomes.

Full model (8 predictors, continuous + categorical): Condition CV AUC = 0.728, Veg response CV AUC = 0.652.
Binary model (5 predictors, simplified binary/categorical — high/low slope, long/short berm, landform, soil development, texture): Condition CV AUC = 0.614, Veg response CV AUC = 0.635.

AUC loss from simplification: Condition delta = 0.113, Veg response delta = 0.016.
Negative permutation importances indicate features that add noise rather than signal after the model has been fit (common for correlated or irrelevant features).}%
  {Feature importance grouped by model type (full / binary), with Condition and Veg-response bars overlaid per predictor. Full model AUC: Condition 0.728, Veg response 0.652; Binary model AUC: 0.614 / 0.635. AUC loss from simplification is small (0.113 / 0.016).}
%% fig_rf_comparison_END


%% fig_rf_performance_START
\figentry%
  {RF Performance Comparison}%
  {fig_rf_performance_comparison.png}%
  {2026-03-27 15:59}%
  {Bar chart comparison of model performance (CV AUC and Holdout AUC) for Full vs Binary random forest models on Intact and Effective outcomes.

Full model (n=8 predictors): Intact CV AUC = 0.728, Effective CV AUC = 0.652
Binary model (n=5 predictors): Intact CV AUC = 0.614, Effective CV AUC = 0.635

Performance difference is minimal (max delta = 0.113 AUC), indicating that simplified binary predictors capture most of the signal. This validates using binary predictors for interpretability without substantial performance loss.}%
  {Model performance comparison showing Full vs Binary RF models achieve similar AUC scores (max delta = 0.113). Binary model provides interpretable results with minimal performance sacrifice.}
%% fig_rf_performance_END


%% fig_rf_full_importance_START
\figentry%
  {RF — feature importance}%
  {fig_rf_full_importance.png}%
  {2026-04-18 09:46}%
  {Random forest model (8 predictors) feature importance for Condition (Intact) and Veg response (Effective) outcomes side by side.
Condition CV AUC = 0.733; top predictors: FA_30_max, Shape_Leng, slope_200.
Veg response CV AUC = 0.655; top predictors: slope_200, Texture, Berm_Length_Class.
Negative importances indicate features that add no signal beyond correlated predictors (effectively zero importance).}%
  {RF (8 predictors) permutation importance: Condition CV AUC 0.733, Veg response CV AUC 0.655. FA_30_max and Shape_Leng dominate for Condition; slope_200 dominates for Veg response.}
%% fig_rf_full_importance_END


%% fig_rf_binary_importance_START
\figentry%
  {RF Binary — feature importance}%
  {fig_rf_binary_importance.png}%
  {2026-03-27 15:59}%
  {Binary random forest model (5 simplified predictors) feature importance for Condition (Intact) and Veg response (Effective) outcomes side by side.
Condition CV AUC = 0.614; top predictors: Long_Berm, Texture, Soil_Development.
Veg response CV AUC = 0.635; top predictors: High_Slope, Texture, Soil_Development.
AUC loss vs full model: Condition 0.113, Veg response 0.016.}%
  {Binary model (5 predictors) permutation importance: Condition CV AUC 0.614, Veg response CV AUC 0.635. Long_Berm and Texture lead for Condition; High_Slope dominates for Veg response.}
%% fig_rf_binary_importance_END


%% FIG_8_START
\figentry%
  {fig8}%
  {fig8_veg_response_by_condition_texture.png}%
  {2026-06-06 11:57}%
  {Panel A: φ = -0.02, p = 0.612. Panel B: Fisher's exact tests per texture (top 3 by sample size); proportions with count labels.}%
  {Structural condition and vegetation response are statistically independent (φ ≈ 0). Sandy loam shows the only significant difference in vegetation response between intact and degraded berms.}
%% FIG_8_END


%% fig_model_comparison_scenarios_START
\figentry%
  {Model comparison — scenario sweep}%
  {fig_model_comparison_scenarios.png}%
  {2026-04-18 10:57}%
  {Dot-plot comparing RF and Logistic L2 across four predictor scenarios for both outcomes.
All-predictors: Condition RF AUC 0.732 ± 0.017, Logistic 0.671 ± 0.032.
All-predictors: Veg response RF AUC 0.653 ± 0.029, Logistic 0.638 ± 0.036.}%
  {RF and Logistic L2 compared across 4 predictor scenarios. All-predictors condition AUC: RF 0.732, Logistic 0.671.}
%% fig_model_comparison_scenarios_END


%% fig_rfecv_feature_selection_START
\figentry%
  {RFECV feature selection effect}%
  {fig_rfecv_feature_selection.png}%
  {2026-03-28 06:19}%
  {Bar chart showing CV AUC before (all 12 predictors) and after RFECV selection for RF and Logistic.
Condition: RFECV selected 12 predictors. Veg response: RFECV selected 4 predictors.}%
  {RFECV reduces predictors (12 → 12 for condition, 4 for veg response) with minimal AUC change.}
%% fig_rfecv_feature_selection_END


%% fig_model_feature_importance_START
\figentry%
  {Feature importance — RF vs Logistic}%
  {fig_model_feature_importance.png}%
  {2026-04-18 10:57}%
  {2×2 feature importance panel: RF permutation importance (top row) and Logistic |coefficient| (bottom row) for Condition and Veg response.
Top RF predictors for Condition: FA_30_max, Shape_Leng, slope_200. Top RF predictors for Veg response: slope_200, Texture, Shape_Leng.}%
  {RF vs Logistic feature importance: top condition predictors are FA_30_max, Shape_Leng, slope_200; top veg response predictors are slope_200, Texture, Shape_Leng.}
%% fig_model_feature_importance_END


%% FIG_4_START
\figentry%
  {fig4}%
  {fig4_pca_biplot.png}%
  {2026-06-06 08:07}%
  {PCA on 8 features (3 numeric + 3 categorical, one-hot encoded, top 3 textures only, n = 630). PC1: 37.5% variance (soil texture–slope axis). PC2: 29.7% variance (landform axis). Cumulative PC1–2: 67.2%.}%
  {PC1 (37.5%) separates berms by soil texture and slope, while PC2 (29.7%) captures landform variation. Intactness shows weak separation along PC1, consistent with soil texture being a secondary predictor of condition. Landform clusters are well separated on PC2 but do not align with intactness, consistent with landform not being independently significant for condition.}
%% FIG_4_END


%% FIG_9_START
\figentry%
  {fig9}%
  {fig9_partial_dependence.png}%
  {2026-04-05 09:27}%
  {Partial dependence plots from Random Forest (n_estimators=300, class_weight='balanced'). Top row: structural condition (Intact probability). Bottom row: vegetation response. Each panel shows the marginal effect of one numeric predictor on the probability of the positive class, averaged over all other features (n = 775).}%
  {Flow accumulation shows a sharp decrease in intactness probability above ~1000 cells, consistent with the optimised FA threshold for condition. Berm length shows decreasing intactness with length. For vegetation response, hillslope gradient shows increasing probability of vegetation response with slope, consistent with steeper slopes enhancing runoff redistribution.}
%% FIG_9_END


%% SI_FA30_POINTPLOT_START
\figentry%
  {Supplementary Figure 1}%
  {SI_fa30_pointplot_by_landform.png}%
  {2026-03-18 15:35}%
  {Point plots show mean $\pm$ 95\% CI of 30\,m flow accumulation (log scale) by landform, grouped by vegetation response (Panel A) and structural condition (Panel B).}%
  {Flood-plain berms experience the highest contributing flow accumulation regardless of outcome, while fan-terrace berms receive the least. Within each landform, effective berms tend to have similar or higher flow accumulation than ineffective berms, and degraded berms tend toward higher values than intact berms, consistent with greater hydraulic stress on degraded structures.}
%% SI_FA30_POINTPLOT_END


%% SI_COUNTPLOTS_START
\figentry%
  {Supplementary Figure 3}%
  {SI_countplots_by_outcome_and_landform.png}%
  {2026-06-06 08:08}%
  {3$\times$3 panel of count plots. Rows: soil development (B horizon vs.\ none), berm length class (short vs.\ long), slope class (shallow vs.\ steep). Columns: vegetation response, structural condition, landform.}%
  {Fan-terrace berms are overwhelmingly found in B-horizon soils, whereas stream-terrace and flood-plain berms dominate in soils without a B horizon. Berm length and slope class are more evenly distributed across outcomes, though short berms are disproportionately intact and steep-slope berms show a higher proportion of vegetation response.}
%% SI_COUNTPLOTS_END


%% FIG_6_START
\figentry%
  {fig6}%
  {fig6_berm_condition.png}%
  {2026-06-06 08:07}%
  {2x3 panel figure. Two-category panels (berm length, slope class, B horizon presence, flow accumulation) use two-sided Fisher's exact tests. Multi-category panels (landform, soil texture) use pairwise two-sided Fisher's exact tests with Benjamini-Hochberg FDR correction. Flow accumulation (30 m) binarized at the fixed 2k cell threshold (two-sided Fisher p = 1.168e-09).}%
  {Six predictors of structural condition (Intact vs Degraded) shown in a 2x3 grid. Row 1: berm length (threshold 60 m), slope class, and landform. Row 2: soil texture (n > 100 per class), B horizon presence, and flow accumulation.}
%% FIG_6_END


%% SI_FA30_LANDFORM_START
\figentry%
  {SI_fa30}%
  {SI_fa30_by_landform.png}%
  {2026-06-06 08:08}%
  {Flow accumulation (30 m, log scale) by landform. Panel A: mean ± 95% CI grouped by vegetation response (point plot). Panel B: mean ± 95% CI grouped by structural condition (bar plot). Flood-plain berms experience the highest contributing flow accumulation regardless of outcome, while fan-terrace berms receive the least.}%
  {Flood-plain berms receive the highest flow accumulation, consistent with their low-gradient valley-floor position. Within each landform, effective berms tend to have similar or higher flow accumulation than ineffective berms, and degraded berms tend toward higher values than intact berms, consistent with greater hydraulic stress on degraded structures.}
%% SI_FA30_LANDFORM_END


%% fig_dropcol_importance_START
\figentry%
  {Drop-column feature importance (RF)}%
  {fig_dropcol_importance.png}%
  {2026-04-18 10:57}%
  {Drop-column feature importance for RF: ΔAUC when each predictor is removed from the full model.
Condition: 12/12 predictors retained (ΔAUC ≥ 0), 0 dropped.
Veg response: 6/12 retained, 6 dropped.}%
  {Drop-column importance: condition retains 12/12 predictors, veg response retains 6/12.}
%% fig_dropcol_importance_END


%% fig_dropcol_feature_selection_START
\figentry%
  {Drop-column feature selection effect}%
  {fig_dropcol_feature_selection.png}%
  {2026-04-18 10:57}%
  {Bar chart showing CV AUC before (all 12 predictors) and after drop-column selection for RF and Logistic.
Condition: kept 12 predictors. Veg response: kept 6 predictors.}%
  {Drop-column selection keeps 12 predictors for condition, 6 for veg response.}
%% fig_dropcol_feature_selection_END


%% FIG_3_START
\figentry%
  {fig3}%
  {fig3_controlled_predictors.png}%
  {2026-07-09 13:36}%
  {Side-by-side comparison of univariate and controlled (LRT) predictor significance for berm structural condition and vegetation response. Panel A: ΔAIC from drop-one likelihood ratio tests (positive = predictor improves model). Panel B: shift from univariate to controlled p-values (circles = univariate, squares = controlled). Six predictors: berm length, slope, landform, soil texture, B-horizon presence, flow accumulation.}%
  {Flow accumulation and berm length are the strongest controlled predictors of condition (largest ΔAIC), while slope drives vegetation response. Several predictors that are significant in univariate tests (e.g. landform, soil texture) lose significance after controlling for confounders, indicating that their apparent effects are partly mediated by correlated site characteristics.}
%% FIG_3_END


%% SI_PREDICTOR_ASSOCIATION_START
\figentry%
  {SI_assoc}%
  {SI_predictor_association.png}%
  {2026-06-10 12:14}%
  {Pairwise association among 8 variables (3 numeric + 3 categorical predictors + 2 binary outcomes — structural condition and vegetation response; top 3 textures only, n = 630). Numeric–numeric pairs use |Spearman ρ|, categorical–categorical pairs use the Bergsma-Wicher bias-corrected Cramer's V, and mixed pairs use the correlation ratio η; all bounded on [0, 1]. Strongest pair: Landform – Soil texture (0.97).}%
  {Spearman's ρ is appropriate only for ordered numeric pairs; for nominal categoricals it reduces to the phi coefficient and discards non-monotonic structure, so Cramer's V and the correlation ratio are used instead. The largest association (Landform–Soil texture, 0.97) indicates predictors are only moderately collinear, while structural condition and vegetation response show only weak association with the predictors and with each other, supporting their joint use in the PCA and predictor-ranking analyses.}
%% SI_PREDICTOR_ASSOCIATION_END


%% SI_THRESHOLD_SENSITIVITY_START
\figentry%
  {SI_threshold_sensitivity}%
  {SI_threshold_sensitivity.png}%
  {2026-06-13 19:24}%
  {Vegetation-response logistic regression refit at SAVI-difference cut-offs from 4% to 12%. Panel A: significance (-log10 p) of slope (Wald) and soil texture (joint LRT) vs. threshold; dashed line p = 0.05, shaded 5-10% reported band, dotted line the 7% cut-off. Panel B: response prevalence vs. threshold (47% at 7%). Across 5-10%, slope is significant (p<0.001) and positive at every cut-off (OR 1.28-1.37); soil texture's joint LRT is significant at [6, 7, 8, 10]%.}%
  {The two headline vegetation-response findings are stable across a band of thresholds: slope is the dominant controlled predictor (significant and positive at every cut-off) and soil texture is second. Texture softens only at the band edges, where the prevalence has drifted furthest from 50% and the joint test loses power, indicating the 7% choice does not drive the conclusions.}
%% SI_THRESHOLD_SENSITIVITY_END

\end{document}
